\newcommand{\irule}[3]{ \inferrule*[right=(#1)]{#2}{#3} }
\def\iaxiom#1#2{\irule{#1}{ }{#2}}

% http://tex.stackexchange.com/questions/3964/mathematical-symbol-for-disjoint-set-union
\makeatletter
\def\moverlay{\mathpalette\mov@rlay}
\def\mov@rlay#1#2{\leavevmode\vtop{%
   \baselineskip\z@skip \lineskiplimit-\maxdimen
   \ialign{\hfil$\m@th#1##$\hfil\cr#2\crcr}}}
\newcommand{\charfusion}[3][\mathord]{
    #1{\ifx#1\mathop\vphantom{#2}\fi
        \mathpalette\mov@rlay{#2\cr#3}
      }
    \ifx#1\mathop\expandafter\displaylimits\fi}
\makeatother

\newcommand{\dcup}{\charfusion[\mathbin]{\cup}{+}}
\newcommand{\bigdcup}{\charfusion[\mathbin]{\bigcup}{+}}

\def\opbin{\boxplus\xspace}
\def\opun{\boxminus\xspace}

% Types
\newcommand{\tInt}{\textsc{Int}\xspace}
\newcommand{\tFloat}{\textsc{Float}\xspace}
\newcommand{\tUnit}{\textsc{Unit}\xspace}
\newcommand{\tBits}{\textsc{Bits}\xspace}
\newcommand{\tChar}{\textsc{Char}\xspace}

\def\structRec#1#2{\textsc{rec}(\&#1) #2}

\newcommand{\tFV}[1]{\textrm{FreeVars}(#1)}

\def\opxor{\mathbin{\char`\^}}

% Constant
\newcommand{\cNil}{\textsc{Nil}}

% Program & language names
\newcommand{\ctonewspeak}{\textsc{c2newspeak}\xspace}
\newcommand{\cpp}{\textsc{cpp}\xspace}
\newcommand{\newspeak}{\textsc{Newspeak}\xspace}

% Newspeak
\newcommand{\npkIf}[3]{\textsc{If}(#1) \{ #2 \} \textsc{Else} \{ #3 \}}
\newcommand{\npkDoWith}[2]{\textsc{Do}\{ #1 \}\textsc{With}~#2:}
\newcommand{\npkWhile}[1]{\textsc{While}(1)\{#1\}}
\newcommand{\npkGoto}[1]{\textsc{Goto}~#1}
\newcommand{\npkReturn}[1]{\textsc{Return}~#1}

% cML
\def\iPass{\textsc{Pass}\xspace}
\def\iIf#1#2#3{\npkIf{#1}{#2}{#3}}
\def\iWhile#1#2{\textsc{While}(#1)\{#2\}}
\def\iReturn#1{\textsc{Return}(#1)}
\def\iDecl#1#2#3{\textsc{Decl}~#1=#2~\textsc{in}\{#3\}}

\def\eNull{\textsc{Null}\xspace}
\def\eUnit{(\,)\xspace}

\def\vRet{\underline{R}\xspace}

% sets pour semantique
\def\sMem{\textsc{Mem}\xspace}
\def\sAddr{\textsc{Addr}\xspace}
\def\sInt{\textsc{Int}\xspace}
\def\sFloat{\textsc{Float}\xspace}
\def\sPhi{\Phi}

% Memory state (instructions)
% #1 : memory
% #2 : instruction
\def\msi#1#2{ \langle #2, #1 \rangle }

% Reduction (expressions)
\def\mm#1#2#3#4{ \msi{#1}{#2} → \msi{#3}{#4} }
\def\mms#1#2{\mm{m}{#1}{m}{#2}}

% Reduction (instructions)
\def\mmi#1#2#3#4{ \mm{#1}{#2}{#3}{#4} }

% Reduction (valeurs gauches)
\def\mmlv#1#2#3#4{ \mm{#1}{#2}{#3}{#4} }

% Typage sém
\def\semtyp#1#2#3{#1 ⊧ #2 : #3}
\def\semtypphi#1#2#3{#1 ⊧_Φ #2 : #3}

% Compatibilité Γ / m
\def\mcomp#1#2{#1 ⊧ #2}

% Reduction contexts
\def\ctxOp#1#2{#1~\opbin~#2}
\def\ctxEmpty{\bullet\xspace}
\def\ctxUnOp#1{\opun~#1}
\def\ctxLvField#1#2{#1.#2}
\def\ctxLvIndex#1#2{#1[#2]}
\def\ctxSet#1#2{#1 ← #2}

% Context substitution
\def\ctxSub#1#2{#1\llparenthesis#2\rrparenthesis}

% Qualifs
\def\qKernel{\textsc{Kernel}\xspace}
\def\qUser{\textsc{User}\xspace}
\def\vTainted#1{\widehat{\diamondsuit}~#1}
\def\eTaint#1{\diamondsuit~#1}

\def\ptrK#1{#1~*}
\def\ptrU#1{#1~@}

% Eval de phrase
\def\ph#1#2#3{#1 ⊩ #2 → #3}

\def\phprog#1#2{⊩ #1 →^* #2}

% typage
\def\ty#1#2#3{#1 ⊢ #2 : #3}
\def\tyi#1#2{#1 ⊢ #2}
\def\typh#1#2#3{#1 ⊢ #2 → #3}

\def\sVal{\textsc{Val}}
\def\sByte{\textsc{Byte}}

\def\powerset#1{\mathcal P \left({#1}\right)}

% Functions with placeholders
\def\phx#1{\ensuremath{#1{\cdot}}}
\def\phxx#1{\ensuremath{#1{\cdot}{\cdot}}}
\def\phxxx#1{\ensuremath{#1{\cdot}{\cdot}{\cdot}}}

% Progs bien formés
\def\WF#1{\textsc{wf}(#1)}
\def\NORET#1{\textsc{NoRet}(#1)}

\newcommand{\insertcode}[1]{\input{gen/#1.pyg.tex}}

% \chapter* + add to toc
%\newcommand{\chapterstar}[1]{\chapter*{#1} \addcontentsline{toc}{chapter}{#1}}

% Register names
\def\defreg#1{\expandafter\def\csname#1\endcsname{\textsc{#1}\xspace}}

\defreg{eax}
\defreg{ebx}
\defreg{ecx}
\defreg{edx}
\defreg{esi}
\defreg{edi}
\defreg{ebp}
\defreg{esp}
\defreg{eip}
\defreg{eflags}

\def\crtrois{\textsc{CR3}}

\newcommand{\ring}[1]{\emph{ring}~#1}

\newcommand{\asminstr}[1]{\texttt{#1}}

\newcommand{\wip}[0]{%
\begin{center}
    \includegraphics[width=5cm]{gen/travaux.pdf}
\end{center}
%\textbf{(wip)}
}

\def\listcons{\texttt{(::)}\xspace}

\tikzstyle{bigbrace} = [decorate,decoration={brace,amplitude=3mm}]

% Debug tikz bounding box
\def\drawbb{\draw[thick] (current bounding box.south west) rectangle (current bounding box.north east);}

% SubFloat + verbatim
\makeatletter
\newbox\sf@box
\newenvironment{SubFloat}[2][]%
{\def\sf@one{#1}%
\def\sf@two{#2}%
\setbox\sf@box\hbox
\bgroup}%
{ \egroup
\ifx\@empty\sf@two\@empty\relax
\def\sf@two{\@empty}
\fi
\ifx\@empty\sf@one\@empty\relax
\subfloat[\sf@two]{\box\sf@box}%
\else
\subfloat[\sf@one][\sf@two]{\box\sf@box}%
\fi}
\makeatother

% {{{ subtitle command
\newcommand\subtitle[1]{%
  \newcommand\thesubtitle{#1}}
% }}}

% Language name
\def\langname{\textsc{Safespeak}\xspace}

% Center vertically + remove header & footer
\newenvironment{headingpage}{
\thispagestyle{simple}
\vspace*{\fill}
}{
\vspace*{\fill}
}

% Hoare triple
\def\hoare#1#2#3{\ensuremath{\{#1\}~#2~\{#3\}}}

\def\eqdef{\overset{\mathrm{def}}{=\joinrel=}\xspace}


% Lens
\def\mkLens#1#2{\langle #1 | #2 \rangle}
\def\setLens#1#2{\textsc{Lens}_{#1, #2}}
\def\lensGetX#1{\ensuremath{\mathrm{get}_{#1}}}
\def\lensPutX#1{\ensuremath{\mathrm{put}_{#1}}}
\def\lensGet#1#2{\lensGetX{#1}(#2)}
\def\lensPut#1#2#3{\lensPutX{#1}(#2, #3)}
\def\lensNodeBig#1#2{
\vcenter{\hbox{
\begin{tikzpicture}[scale=0.7]
\draw[#1] (-0.5,-0.5) rectangle ++ (1,1);
\draw[#2] (0,0) circle (0.3);
\end{tikzpicture}
}}
}
\def\lensInner#1{
\vcenter{\hbox{
\begin{tikzpicture}[scale=0.7]
\draw[#1] (0, 0) circle (0.3);
\end{tikzpicture}
}}
}
\def\lensNodeBigIdx#1#2#3#4#5{
\vcenter{\hbox{
\begin{tikzpicture}[scale=0.7]
  \draw[#1] (-1.25,-1.25) rectangle ++ (2.5,2.5);
  \node[#2, draw, star, minimum size=6mm] at (-0.6, 0.6) {};
  \node[#3, draw, star, minimum size=6mm] at ( 0.6, 0.6) {};
  \node[#4, draw, circle, minimum size=6mm] at (-0.6,-0.6) {};
  \node[#5, draw, regular polygon, regular polygon sides=5, minimum size=6mm] at ( 0.6,-0.6) {};
  \node at (-2,  0.6) {a};
  \node at ( 2,  0.6) {b};
  \node at (-2, -0.6) {c};
  \node at ( 2, -0.6) {d};
\end{tikzpicture}
}}
}
\def\lensInnerStar#1{
\vcenter{\hbox{
\begin{tikzpicture}[scale=0.7]
\node[draw, #1, star] {};
\end{tikzpicture}
}}
}


% Erreurs

\def\serr#1{
  \switch{#1}
    {sec}{Ω_{sec}}
    {array}{Ω_{array}}
    {ptr}{Ω_{ptr}}
    {div}{Ω_{div}}
    {field}{Ω_{field}}
    {var}{Ω_{var}}
    {typ}{Ω_{typ}}
    {END}
    % important : étendre la description dans le chapitre 4 et la figure de
    % liste
}

%\def\uGet#1#2{#1 ⇐_U #2}
%\def\uPut#1#2{#1 ⇒_U #2}
\def\cfuF{\mathrm{copy\_from\_user}\xspace}
\def\ctuF{\mathrm{copy\_to\_user}\xspace}
\def\uGet#1#2{\cfuF(#1, #2)}
\def\uPut#1#2{\ctuF(#1, #2)}

% Litteral notation

\def\eStruct#1{\{#1\}}
\def\eArray#1{[#1]}
\def\tStruct#1{\eStruct{#1}}
\def\eFun#1#2{\mathrm{fun} (#1) \{ #2 \}}

\def\gramorx{\hspace{3mm}\gramor\xspace}
\def\stFun#1{\textsc{Fun}_n}
\def\tComp#1#2{#1 \triangleright #2}

% Box on top on rules fig
\def\ruleheader#1{
  \begin{minipage}{\linewidth}
    \raggedleft
    \fbox{#1}
  \end{minipage}
}

\def\termFst#1{\ensuremath{\textrm{fst}~#1}}
\def\termSnd#1{\ensuremath{\textrm{snd}~#1}}

\def\ptrtype{\texttt{ptrtype}\xspace}

\def\addrInt{\textrm{addr}^\natural \xspace}
\def\interpPtr#1{\textrm{interp}_{(#1)}}

\tikzset{>=latex}

\def\partccl#1{
  \cleartorecto
  \chapter*{Conclusion de la partie~\ref{#1}}
  \addcontentsline{toc}{chapter}{Conclusion de la partie~\ref{#1}}
}

\def\hexa#1{\texttt{0x#1}}
\def\floatop#1{#1.}

\def\cLookup#1#2{\mathrm{Lookup} (#1, #2)}
\def\cExtend#1#2#3{\mathrm{Extend}(#1, #2 ↦ #3)}
\def\cCleanup#1{\mathrm{Cleanup}(#1)}

% {{{
\def\jolibreak{
\fancybreak{\decosix \quad \decosix \quad \decosix}
}
% }}}


\def\cDom#1{\mathrm{dom}(#1)}

\def\cVisible#1{\mathrm{Visible}(#1)}

\def\cRepr#1{\mathrm{Repr}(#1)}

\def\cPush#1#2{\mathrm{Push}(#1, #2)}

\def\cPop#1{\mathrm{Pop}(#1)}

\def\cH#1{\mathcal{H}(#1)}
\def\cHphi#1{\mathcal{H}_Φ(#1)}
\def\cV#1#2{\mathrm{CleanV}_{#1}(#2)}
\def\cP#1#2{\mathrm{CleanP}(#1, #2)}
\def\cLP#1#2{\mathrm{CleanLP}(#1, #2)}
\def\cVP#1#2{\mathrm{CleanVP}(#1, #2)}

\def\cCleanVar#1#2{\mathrm{CleanVar}(#1, #2)}
\def\cVV#1#2{\mathrm{CleanVarV}(#1, #2)}

\def\wphi#1{\widehat{\&}~#1}

\def\extendLocal#1#2#3{#1, \mathrm{local}~#2:#3}
\def\extendGlobal#1#2#3{#1, \mathrm{global}~#2:#3}

\def\cVar#1{\mathrm{Adr}(#1)}

\newcommand*\circled[1]{\tikz[baseline=(char.base)]{
            \node[shape=circle,draw,inner sep=2pt] (char) {#1};}}

\newcommand{\sList}[1]{#1^*}
